\section{General Introduction to Monte Carlo Methods}
Our goal is to evaluate the expectation value of an observable $O[U]$ in the quantized Euclidian gauge field theory on a lattice. This is written as a path integral
\begin{equation}
    \braket{O} = \frac{1}{Z} \int \mathcal{D}[U] \; \exp \left( -S[U] \right) O[U], \label{eq: path_integral}
\end{equation}
where the partition function is defined as
\begin{equation}
    Z = \int \mathcal{D}[U] \exp \left( -S[U] \right). \label{eq: partition_function}
\end{equation}
Here, $S[U]$ denotes the Euclidian action and $\mathcal{D}[U]$ represents the functional integration over all field configurations. In the general case, these integrals cannot be solved analytically and depending on the application, these can be high-dimensional \cite{GattingerLang}. \\\\
Using probability theory, we can approximate the integral of a function $f(x)$ by averaging the functional values $f(x_n)$ at values $x_n$, which are randomly chosen according to a uniform distribution 
\begin{equation}
    \rho_u(x_n) = \frac{1}{b-a}. 
\end{equation}
The integral is approximated as: 
\begin{equation}
    \frac{1}{b-a}\int_a^b dx \; f(x) = \lim_{N \rightarrow \infty} \frac{1}{N}\sum_{n=1}^N f(x_n). 
\end{equation}
The exact mean is replaced by the sample mean, the calculation is restricted to a subsample and a finite number $N$. By the Strong Law of Large Numbers, the Monte Carlo estimator
\[
\frac{1}{N} \sum_{n=1}^{N} f(x_n)
\]
converges almost surely to the exact expectation value as $N \to \infty$, and the Central Limit Theorem implies an error scaling of $\mathcal{O}\!\left( \frac{1}{\sqrt{N}} \right)$. This means to increase the accuracy by a factor of two, one has to use four times as many samples \cite{GattingerLang}. 
\subsection{Example of a Simple Monte Carlo Integral}
To illustrate the basic idea of Monte Carlo integration, we consider a simple one-dimensional example. Suppose we want to estimate the integral
\begin{equation}
\int_0^1 \exp(-x^2) \, dx.
\end{equation}
We restrict the upper limit to $1$ to match the support of a uniform random variable. Let $X$ be a random variable uniformly distributed on the interval $(0,1)$. Its probability density function (pdf) is given by
\begin{equation}
f(x) = 
\begin{cases}
1 & \text{if } x \in (0,1), \\
0 & \text{otherwise}.
\end{cases}
\end{equation}
Then the integral can be written as an expectation, 
\begin{equation}
\int_0^1 \exp(-x^2) \, dx = \mathbb{E}[\exp(-X^2)].
\end{equation}
A Monte Carlo estimator using $N$ independent samples $X_1, X_2, \dots, X_N$ from the uniform distribution is
\begin{equation}
I_N = \frac{1}{N} \sum_{i=1}^N \exp(-X_i^2).
\end{equation}
By the law of large numbers, this estimator converges to the true value as $N \to \infty$, 
\[
I_N \xrightarrow[]{N \to \infty} \int_0^1 \exp(-x^2) \, dx.
\]
Moreover, by the central limit theorem, the statistical error scales as $\mathcal{O}\left( \frac{1}{\sqrt{N}}\right)$, which is independent of the dimensionality of the integral. The not scaling with dimension is an important reason for the use of Monte Carlo methods in lattice field theory. This simple example illustrates how Monte Carlo methods approximate integrals by replacing them with sample averages \cite{Monte_carlo}.
\subsection{Importance Sampling}
In our path integral in Equation \ref{eq: path_integral}, we need to account for the Boltzmann factor $\exp(-S)$. This is a weighting factor, depending on the action $S$, of the different field configurations. 
The space of all configurations is not uniformely explored, but according to the weights determined by the action $S$. The main idea of importance sampling is to approximate a large sum with a comparatively small amount of configurations, sampled according to this weight factor. \\\\
For a general probability distribution $\rho(x)$, the expectation value of a function $f(x)$ is
\begin{equation}
    \braket{f}_\rho = \frac{\int_a^b dx \; \rho(x) f(x)}{\int_a^b dx \; f(x)}. 
\end{equation}
In the importance sampling Monte Carlo integration, this is approximated as follows: 
\begin{equation}
    \braket{f}_\rho = \lim_{N \rightarrow \infty} \frac{1}{N} \sum_{n = 1}^N f(x_n), 
\end{equation}
where $x \in (a, b)$ sampled according to the normalized probability density $\frac{\rho(x) dx}{\int_a^b dx \rho(x)}$. Our path integral (Equation \ref{eq: path_integral}) is of this form, so we can apply this method: 
\begin{equation}
    \braket{O} = \lim_{N \rightarrow \infty} \frac{1}{N} \sum_{n=1}^N O[U_n], 
\end{equation}
with $U_n$ sampled according to the probability distribution density 
\begin{equation}
    dP(U) = \frac{e^{-S[U]} \mathcal{D}[U]}{\int \mathcal{D}[U] e^{-S[U]}}, 
\end{equation}
which is called the Gibbs-measure \cite{GattingerLang}. 
\subsection{Markov Chains}
The configurations $U_n$ are generated in a sequence, a so-called homogeneous Markov chain: $U_0 \rightarrow U_1 \rightarrow U_2 \rightarrow \dots $. 
The first configuration $U_0$ is chosen randomly, the update from one configurations to the next is called the Markov step. The next configuration only depends on the present one, not one
the history of the process. This condition is written as 
\begin{equation}
    P(U_n = U'| U_{n-1} = U) = T(U'|U). 
\end{equation}
The configuration sequence explores the space of all configurations according to the probability density defined by the action. Eventually, the Markov chain will reach an equilibrium, after that, it cannot have sinks or sources of probability. This gives rise to a balancing condition: 
\begin{equation}
    \sum_U T(U|U') P(U) = \sum_U T(U'|U) P(U'). \label{eq: balancing_condition}
\end{equation}
A solution for the balancing condition in Equation \ref{eq: balancing_condition} can be obtained by requiring term-wise equality: 
\begin{equation}
    T(U|U') P(U) = T(U'|U) P(U'). 
\end{equation}
This is called detailed balance condition, it is the most commonly used solution. \\
Using the normalization property, we find in addition that 
\begin{equation}
    P(U') = \sum_U T(U|U') P(U). 
\end{equation}
The equilibrium distribution $P(U)$ is a fixed point of the Markov process, once the equilibrium is reached, the system stays there forever. \\
In order to give correct result, the Markov chain must be aperiodic, i.e. be able to reach all configurations in the configuration space. This is the case iff the transition matrix $T(U|U')$ is positive-semidefinite. This property is called strong ergodicity \cite{GattingerLang}. 
\section{The Metropolis Algorithm}
The Metropolis Algorithm is used to obtain a new configuration $U_{n+1}$ from the current configuration $U_n$. We use $P[U] \propto \exp \left( -S[U]\right)$. The steps are as follows: 
\begin{enumerate}
    \item Choose a candidate configuration $U'$ according to some a priori selection probability $T_0(U'|U)$, where $U = U_n$. 
    \item Accept the candidate configuration $U'$ as the new configuration $U_{n+1}$ with the following probability: 
    \begin{equation}
        T_A(U'|U) = \min \left( 1, \frac{T_0(U|U') \exp( - S[U'])}{T_0(U'|U)\exp (-S[U])}\right). 
    \end{equation}
    If the change is rejected, keep the old configuration, set $U_{n+1} = U_n$. 
    \item Repeat. 
\end{enumerate}
The detailed balance condition is fulfilled: 
\begin{align}
    T(U'|U) \exp(-S[U]) &= T_0(U'|U) \min \left( 1, \frac{T_0(U|U') \exp( - S[U'])}{T_0(U'|U)\exp (-S[U])}\right) \exp (-S[U]) 
    \\ &= \min \left( T_0(U|U') \exp( - S[U']), T_0(U'|U)\exp (-S[U]) \right) \\ &= T(U|U') \exp(-S[U'])
\end{align}
We have used the symmetry of min and the positivity of all factors. \\\\
Often, we have $T_0(U|U') = T_0(U'|U)$. In this case, the transition probability simplifies to
\begin{equation}
    T_A(U'|U) = \min \left( 1, \exp (-\Delta S) \right), \quad \text{with} \quad \Delta S = S[U'] - S[U].
\end{equation}
The decision to accept or reject only depends on the change in the action between the two configurations. Any normalization constant in $P[U]$ will cancel automatically, note that this is crucial in field theory where the partition function $Z$ is unknown. 
In practice, the acceptance rate is strongly dependent on the proposed change. Very large proposed updates are often rejected, very small ones lead to long autocorrelation times. Hence we often aim for acceptance rates between 50\% - 80\%. 
\\\\
A warm-up period (thermalization) is required to eliminate dependence on the initial configuration. After that, configurations are still correlated, and one must estimate autocorrelation times when computing errors.
\\\\
It can be shown that the Metropolis algorithm fulfills ergodicity, meaning that any configuration is reachable from any other configuration with positive probability in a finite number of steps. Because the transition kernel satisfies detailed balance with respect to $\exp (-S[U])$, this probability distribution is the stationary distribution of the Markov chain. Together with ergodicity, this ensures convergence of the algorithm to the desired ensemble. \cite{GattingerLang}